%%%%%%%%%%%%%%%%%%%%%%%%%%%%%%%%%%%%%%%%%%%%%%%%%%%%%%%%%%%%%%%%%%%%%
\section{Historical Development}
%%%%%%%%%%%%%%%%%%%%%%%%%%%%%%%%%%%%%%%%%%%%%%%%%%%%%%%%%%%%%%%%%%%%%

The possibility that quantum mechanics may arise from an underlying
stochastic process has attracted continuous attention since the
development of Brownian motion and modern probability theory. Although
the mathematical foundations of stochastic calculus were established by
Wiener, Itô and Stratonovich, the first systematic attempt to derive the
Schrödinger equation from stochastic dynamics was made by Edward Nelson
in the 1960s.

Nelson's construction demonstrated that nonrelativistic quantum
mechanics may be formulated as a conservative diffusion process on
Euclidean configuration space. In this formulation, particles follow
continuous but nondifferentiable trajectories, and the wave function
emerges from the probability density together with the current velocity.

The success of the nonrelativistic theory immediately suggested the
possibility of extending stochastic mechanics to relativistic systems.
This extension proved to be considerably more difficult than originally
anticipated. The principal obstacle is that Brownian motion is naturally
defined on spaces endowed with a positive-definite metric, whereas
Minkowski space possesses an indefinite Lorentzian metric.

Over the past six decades numerous approaches have been proposed,
including proper-time formulations, Lorentz-invariant diffusion
processes, stochastic differential geometry on Lorentzian manifolds,
relativistic kinetic theory, and off-shell formulations based upon the
Stueckelberg--Horwitz--Piron framework.

Despite substantial progress, no formulation has yet achieved the same
degree of completeness and acceptance as Nelson's original
nonrelativistic theory. The purpose of this chapter is to review these
developments, identify their strengths and limitations, and establish
the foundations upon which the SHP stochastic theory developed in the
following chapters will be constructed.



%%%%%%%%%%%%%%%%%%%%%%%%%%%%%%%%%%%%%%%%%%%%%%%%%%%%%%%%%%%%%%%%%%%%%%%%%%%%%%%%
\section{Wiener Processes and the Mathematical Foundations of Diffusion}
%%%%%%%%%%%%%%%%%%%%%%%%%%%%%%%%%%%%%%%%%%%%%%%%%%%%%%%%%%%%%%%%%%%%%%%%%%%%%%%%

The mathematical foundation of stochastic mechanics is the Wiener
process. Before discussing relativistic stochastic dynamics, it is
therefore necessary to review the essential properties of Brownian
motion and the associated stochastic calculus. Throughout this chapter
we emphasize the physical interpretation of the mathematical
construction, while maintaining sufficient rigor for later applications.

%%%%%%%%%%%%%%%%%%%%%%%%%%%%%%%%%%%%%%%%%%%%%%%%%%%%%%%%%%%%%%%%%%%%%%%%%%%%%%%%

\subsection{Brownian Motion}

The phenomenon of Brownian motion was first observed by Robert Brown in
1827 during microscopic studies of pollen grains suspended in water.
Although initially regarded as a biological phenomenon, it was later
recognized as a universal consequence of molecular agitation.

The first quantitative description was given independently by Einstein
(1905) and Smoluchowski (1906), who demonstrated that the observed
motion follows from the random impacts of surrounding molecules.
Einstein's analysis established the connection between microscopic
kinetics and macroscopic diffusion and provided one of the earliest
quantitative confirmations of the atomic theory of matter.

The modern mathematical formulation was developed by Wiener, who
introduced a probability measure on the space of continuous paths. The
resulting stochastic process is now known as the Wiener process.

%%%%%%%%%%%%%%%%%%%%%%%%%%%%%%%%%%%%%%%%%%%%%%%%%%%%%%%%%%%%%%%%%%%%%%%%%%%%%%%%

\subsection{Definition of the Wiener Process}

A real stochastic process

\[
W(t),
\qquad t\ge0,
\]

is called a Wiener process if it satisfies the following properties.

\begin{enumerate}

\item
The process begins at the origin,

\[
W(0)=0.
\]

\item
The increments are independent,

\[
W(t_2)-W(t_1)
\]

is statistically independent of

\[
W(t_4)-W(t_3)
\]

whenever

\[
t_1<t_2<t_3<t_4.
\]

\item
The increments are stationary,

\[
W(t+\Delta t)-W(t)
\]

depends only upon

\[
\Delta t.
\]

\item
Each increment possesses a Gaussian distribution,

\[
W(t+\Delta t)-W(t)
\sim
{\cal N}(0,\Delta t).
\]

\item
With probability one, the trajectories are continuous but nowhere
differentiable.

\end{enumerate}

These properties completely characterize Brownian motion.

%%%%%%%%%%%%%%%%%%%%%%%%%%%%%%%%%%%%%%%%%%%%%%%%%%%%%%%%%%%%%%%%%%%%%%%%%%%%%%%%

\subsection{Gaussian Increments}

The probability density for a finite increment is

\[
P(\Delta W)
=
\frac{1}
{\sqrt{2\pi\Delta t}}
\exp
\left(
-
\frac{(\Delta W)^2}
{2\Delta t}
\right).
\]

Consequently,

\[
\langle\Delta W\rangle=0,
\]

and

\[
\boxed{
\langle(\Delta W)^2\rangle
=
\Delta t.
}
\]

Higher moments satisfy

\[
\langle(\Delta W)^{2n+1}\rangle=0,
\]

while

\[
\langle(\Delta W)^4\rangle
=
3(\Delta t)^2.
\]

These relations play a central role in the derivation of Itô's lemma.

%%%%%%%%%%%%%%%%%%%%%%%%%%%%%%%%%%%%%%%%%%%%%%%%%%%%%%%%%%%%%%%%%%%%%%%%%%%%%%%%

\subsection{Quadratic Variation}

Unlike ordinary differentiable functions,

Brownian trajectories possess finite quadratic variation.

Indeed,

\[
(dW)^2
=
dt,
\]

while

\[
(dW)^3
=
0,
\]

and

\[
(dW)^4
=
O(dt^2).
\]

The identity

\[
\boxed{
(dW)^2=dt
}
\]

is the defining algebraic property of Itô calculus.

It distinguishes stochastic analysis from ordinary differential
calculus and is responsible for the appearance of second-order
derivative terms in stochastic differential equations.

%%%%%%%%%%%%%%%%%%%%%%%%%%%%%%%%%%%%%%%%%%%%%%%%%%%%%%%%%%%%%%%%%%%%%%%%%%%%%%%%

\subsection{Diffusion Constant}

Physical Brownian motion is usually written in the form

\[
dX
=
b(X,t)\,dt
+
\sqrt{2D}\,dW,
\]

where

\[
b(X,t)
\]

is the drift velocity and

\[
D
\]

is the diffusion coefficient.

The statistical moments become

\[
\langle dX\rangle
=
b\,dt,
\]

and

\[
\boxed{
\langle(dX)^2\rangle
=
2D\,dt.
}
\]

The parameter

\[
D
\]

controls the strength of stochastic fluctuations.

In Nelson's stochastic mechanics,

it assumes the distinguished value

\[
\boxed{
D=\frac{\hbar}{2m},
}
\]

thereby establishing the connection between stochastic diffusion and
quantum mechanics.

%%%%%%%%%%%%%%%%%%%%%%%%%%%%%%%%%%%%%%%%%%%%%%%%%%%%%%%%%%%%%%%%%%%%%%%%%%%%%%%%

\subsection{Markov Property}

The Wiener process is Markovian.

The future evolution depends only upon the present state,

not upon the previous history.

Consequently,

\[
P(x_3,t_3|x_2,t_2;x_1,t_1)
=
P(x_3,t_3|x_2,t_2).
\]

The corresponding transition probability satisfies the
Chapman--Kolmogorov equation,

\[
P(x_3,t_3|x_1,t_1)
=
\int
P(x_3,t_3|x_2,t_2)
P(x_2,t_2|x_1,t_1)
dx_2.
\]

This property forms the probabilistic basis of the Fokker--Planck
equation.



%%%%%%%%%%%%%%%%%%%%%%%%%%%%%%%%%%%%%%%%%%%%%%%%%%%%%%%%%%%%%%%%%%%%%%
\chapter{Relativistic Stochastic Mechanics}
%%%%%%%%%%%%%%%%%%%%%%%%%%%%%%%%%%%%%%%%%%%%%%%%%%%%%%%%%%%%%%%%%%%%%%

\section{Why Stochastic Mechanics?}

Classical mechanics assumes that the state of a particle can be
specified exactly by its position and momentum at every instant of
time. Once the initial conditions are given, Hamilton's equations
determine the future motion uniquely.

This deterministic picture has proved extraordinarily successful.
Nevertheless, it raises an obvious question. Is the deterministic
trajectory a fundamental property of nature, or is it merely an
effective description obtained after averaging over microscopic degrees
of freedom that remain hidden?

Brownian motion provides a familiar example in which deterministic
mechanics ceases to be adequate. A microscopic particle suspended in a
fluid is continually bombarded by an enormous number of surrounding
molecules. Although every individual collision obeys Newton's laws, the
combined effect of billions of collisions per second is effectively
random. The detailed microscopic dynamics becomes inaccessible, while
its statistical properties remain well defined.

The natural question is therefore whether quantum phenomena might admit
a similar description. Instead of postulating a wave function from the
outset, one may ask whether quantum mechanics could emerge from an
underlying stochastic dynamics.

This idea has a long history and reached its most complete formulation
in the stochastic mechanics developed by Nelson. In that theory the
particle follows a continuous but non-differentiable trajectory, while
the Schrödinger equation appears as the equation governing the
probability distribution of an ensemble of such trajectories.

The success of Nelson's theory immediately suggests extending the same
program to relativistic systems. Unfortunately, this turns out to be
far from straightforward. The difficulties are not merely technical but
reflect fundamental differences between Euclidean probability theory and
the geometry of Minkowski space.

The purpose of this chapter is to examine these difficulties and to
review the principal attempts to overcome them. Only after identifying
the limitations of existing approaches shall we construct a stochastic
theory within the SHP framework.



%%%%%%%%%%%%%%%%%%%%%%%%%%%%%%%%%%%%%%%%%%%%%%%%%%%%%%%%%%%%%%%%%%%%%%%%%%
\section{What is Fluctuating?}
%%%%%%%%%%%%%%%%%%%%%%%%%%%%%%%%%%%%%%%%%%%%%%%%%%%%%%%%%%%%%%%%%%%%%%%%%%

Every stochastic theory begins with an assumption that some physical
quantity is subject to random fluctuations. The choice of this quantity
is not unique and determines the structure of the entire theory.

In ordinary Brownian motion the answer is well known. The particle is
continually struck by surrounding molecules. The microscopic collisions
produce a rapidly fluctuating force, while the observed random motion of
the particle is only the accumulated effect of these countless impacts.

Thus the position is not fundamentally random. It becomes random because
the force acting on the particle is random.

The same question must be asked in any relativistic theory.

What is the physical quantity that fluctuates?


Suppose that one assumes the space-time position itself to undergo a
Brownian motion.

The world line then becomes a continuous but nowhere differentiable
curve. This is the viewpoint adopted in Nelson's original stochastic
mechanics.

Its principal advantage is simplicity. The probability density is
defined directly in configuration space, and the Schrödinger equation
emerges naturally from the dynamics of the ensemble.

The difficulty appears only after relativity is introduced.

A Brownian process requires a positive-definite covariance matrix,
whereas Minkowski space possesses an indefinite metric. Consequently the
construction that works perfectly in Euclidean space cannot be
transferred directly to space-time.


Instead of assuming that the position fluctuates, one may regard the
velocity as the stochastic variable.

The position then remains the integral of the random velocity.

This viewpoint has an important physical motivation.

Random collisions change the particle's velocity directly, while the
position changes only as a consequence of the altered velocity.

Several relativistic diffusion theories are based upon this idea.

The principal difficulty is that relativistic velocities cannot vary
arbitrarily. They are constrained by the geometry of space-time and by
the requirement that no trajectory exceed the speed of light.

The stochastic process must therefore evolve on the manifold of allowed
four-velocities rather than in an unconstrained Euclidean space.


From the viewpoint of Hamiltonian mechanics, momentum appears to be the
most natural quantity to fluctuate.

Microscopic collisions exchange momentum between the particle and its
environment.

Hamilton's equations then determine the corresponding evolution of the
trajectory.

This viewpoint is particularly attractive because it preserves the
Hamiltonian structure of the dynamics.

Relativistic kinetic theory and Langevin equations are usually
constructed in this manner.

The price paid for this simplicity is that the momentum cannot vary
independently of the relativistic mass-shell constraint.

One must therefore decide whether the stochastic dynamics is confined to
the mass shell or whether off-shell motion is permitted.


An even more physical description regards neither position nor momentum
as fundamentally stochastic.

Instead, the fluctuating quantity is the force itself.

This picture corresponds most closely to the microscopic origin of
Brownian motion.

Each collision produces a short impulse.

The observable force is therefore the sum of an enormous number of
rapidly fluctuating contributions.

Position and momentum become stochastic only after integration over
time.

The resulting equations are known as Langevin equations.


The SHP framework suggests yet another possibility.

The Hamiltonian formulation already introduces an invariant evolution
parameter together with a complete phase-space description.

Instead of modifying the geometric structure of space-time, one may
therefore regard the Hamiltonian evolution itself as being driven by
microscopic fluctuations.

The stochastic process then develops in the invariant parameter
\(\tau\), while the four-dimensional covariance of the theory remains
manifest.

This viewpoint combines naturally with the Hamiltonian structure
developed in the previous chapters and will serve as the foundation for
the stochastic theory developed later in this monograph.


%%%%%%%%%%%%%%%%%%%%%%%%%%%%%%%%%%%%%%%%%%%%%%%%%%%%%%%%%%%%%%%%%%%%%%%%%%%%%%%%
\section{Why Relativistic Brownian Motion is Difficult}
%%%%%%%%%%%%%%%%%%%%%%%%%%%%%%%%%%%%%%%%%%%%%%%%%%%%%%%%%%%%%%%%%%%%%%%%%%%%%%%%

The extension of Brownian motion from Newtonian mechanics to special
relativity appears, at first sight, to be straightforward. One might
expect that every three-dimensional quantity should simply be replaced
by its four-dimensional counterpart.

Experience has shown that this expectation is incorrect.

The difficulties encountered by relativistic stochastic mechanics are
not technical details of probability theory but arise from the
fundamental structure of relativistic space-time itself.

In this section we discuss the physical origin of these difficulties
before reviewing the various mathematical approaches that have been
proposed to overcome them.


%%%%%%%%%%%%%%%%%%%%%%%%%%%%%%%%%%%%%%%%%%%%%%%%%%%%%%%%%%%%%%%%%%%%%%%%%%%%%

\subsection{Which Time?}

Every stochastic process describes the evolution of a system with
respect to some parameter.

In ordinary Brownian motion that parameter is simply the Newtonian time

\[
t.
\]

This choice is so natural that it is rarely questioned.

Special relativity changes the situation completely.

Different inertial observers assign different values to the coordinate
time.

Consequently there is no unique universal time with respect to which the
random process may evolve.

One may attempt to use the proper time of the particle.

However, proper time is defined only after the trajectory is known,
whereas the trajectory itself is the unknown stochastic process.

Thus the evolution parameter cannot be specified independently of the
motion.

This circularity lies at the heart of many relativistic stochastic
formulations.


The SHP framework provides a remarkably simple resolution of this
difficulty.

Instead of using either the coordinate time or the proper time, the
theory introduces an invariant evolution parameter

\[
\tau ,
\]

which exists independently of the particle trajectory.

The stochastic process therefore evolves with respect to a parameter
that is both Lorentz invariant and known before the trajectory is
constructed.

One of the principal motivations for combining stochastic mechanics with
the SHP formalism is precisely the existence of this invariant
evolution parameter.


Ordinary relativistic mechanics describes particles whose invariant mass
is fixed.

Random interactions, however, continually exchange energy and momentum
with the surrounding environment.

One must therefore answer an important question.

Does the particle remain on the same mass shell after every microscopic
collision?

If the answer is yes, an explicit constraint must accompany the
stochastic dynamics.

If the answer is no, the theory must allow the particle to evolve
off-shell.

Different relativistic stochastic theories make different choices at
this point.

The SHP formalism adopts the second viewpoint from the outset.


Probability theory requires the covariance of the random increments to
be positive.

This condition guarantees that the mean-square fluctuation is always
non-negative.

Minkowski space possesses an indefinite metric.

Consequently one cannot simply replace

\[
\delta_{ij}
\]

by

\[
\eta_{\mu\nu}
\]

in the covariance matrix of a Wiener process.

This is not merely a technical inconvenience.

It reflects the fundamental difference between Euclidean geometry and
Lorentzian geometry.

Nearly every relativistic stochastic theory must confront this problem
in one form or another.


The preceding discussion demonstrates that relativistic stochastic
mechanics is not a single problem but a collection of interconnected
problems.

One must choose

the evolution parameter,

the stochastic variables,

the underlying geometric manifold,

the treatment of the mass shell,

and the relation between the stochastic process and quantum dynamics.

Different formulations make different choices and therefore answer
different physical questions.

The following sections review the principal approaches that have been
developed during the past six decades.




